\documentclass[14pt]{extarticle}
\def\activityid{F02-M-v4}\usepackage[letterpaper,margin=.65in,top=.60in,bottom=.65in]{geometry}
\usepackage{fix-cm}
\usepackage[T1]{fontenc}\usepackage[default]{sourcesanspro}
\usepackage{microtype,tikz,array,tabularx,fancyhdr,amsmath,amssymb}
\usepackage[hidelinks]{hyperref}\usetikzlibrary{calc,arrows.meta}
\setlength{\parindent}{0pt}\setlength{\parskip}{7pt}\setlength{\headheight}{0pt}\setlength{\footskip}{26pt}
\setlength{\emergencystretch}{2em}\pagestyle{fancy}\fancyhf{}\renewcommand{\headrulewidth}{0pt}\renewcommand{\footrulewidth}{.4pt}
\fancyfoot[L]{\scriptsize Bellingham Math Circle / Week 2 / \activityid}\fancyfoot[R]{\scriptsize \thepage}
\definecolor{ink}{gray}{.12}\definecolor{soft}{gray}{.95}\color{ink}
\newcommand{\PageTitle}[2]{{\fontsize{10}{12}\selectfont\bfseries WEEK 2\quad TOGGLE LAB\hfill #1\par}\vspace{3pt}{\fontsize{26}{29}\selectfont\bfseries #2\par}\vspace{2pt}}
\newcommand{\NameLine}{{\small Name: \rule{2.65in}{.4pt}\hfill Date: \rule{1.35in}{.4pt}\par}\vspace{4pt}}
\newcommand{\Task}[2]{\par\vspace{3pt}{\large\bfseries #1}\par #2\par}
\newcommand{\Subhead}[1]{\par\vspace{3pt}{\large\bfseries #1}\par}
\newcommand{\RuleBox}[1]{\par\begingroup\setlength{\fboxsep}{8pt}\colorbox{soft}{\parbox{\dimexpr\linewidth-16pt\relax}{#1}}\endgroup\par}
\newcommand{\WriteBox}[2]{\par\textbf{#1}\par\vspace{2pt}\begin{tikzpicture}\draw[black!40,line width=.5pt] (0,0) rectangle (\dimexpr\linewidth-.5pt\relax,#2);\end{tikzpicture}\par}
\newcommand{\StudentFooter}{\vfill{\small Try it with objects. Explain by pointing, talking, or drawing.}\par}
\newcommand{\LampRing}[3]{\begin{tikzpicture}[x=1in,y=1in]
\foreach \i in {1,...,#1}{\coordinate (v\i) at ({90-360*(\i-1)/#1}:#2);}
\foreach \i in {1,...,#1}{\pgfmathtruncatemacro{\j}{mod(\i,#1)+1}\draw[thick] (v\i)--node[midway,fill=white,inner sep=2pt,font=\small]{\i\j}(v\j);}
\foreach \i in {1,...,#1}{\node[circle,draw,fill=white,minimum size=.52in,inner sep=0pt,font=\bfseries] at(v\i){\i};}
\foreach \i in {#3}{\node[circle,draw,fill=ink,text=white,minimum size=.52in,inner sep=0pt,font=\bfseries] at(v\i){\i};}
\end{tikzpicture}}
\newcommand{\Lights}[1]{\begin{tikzpicture}[x=.55in,y=.55in,baseline=-.10in]\foreach \v [count=\i] in {#1}{\ifnum\v=1\fill (\i,0) circle(.16in);\else\draw (\i,0) circle(.16in);\fi}\end{tikzpicture}}
\newcommand{\ClockRing}[2]{\begin{tikzpicture}[x=1in,y=1in]\draw[black!30] (0,0) circle(#2);\pgfmathtruncatemacro{\n}{#1-1}\foreach\i in {0,...,\n}{\node[circle,draw,fill=white,minimum size=.35in,inner sep=0pt] at({90-360*\i/#1}:#2){\i};}\draw[-{Stealth},thick] (-.35,.15) arc(155:25:.38);\node[font=\small] at(0,-.18){clockwise};\end{tikzpicture}}
\newcommand{\Key}[2]{\begin{center}\renewcommand{\arraystretch}{1.55}\begin{tabular}{c|*{#1}{c}}in&#2\end{tabular}\end{center}}


% v3 student pages use one running header and numbered tasks only.
\newcommand{\StudentHeader}[1]{{\fontsize{10}{12}\selectfont\bfseries Week 2 / Toggle lab / #1\par}\vspace{10pt}}
\newcommand{\Problem}[2]{\par\vspace{4pt}\textbf{Problem #1:} #2\par}
\newcommand{\Space}[1]{\begin{tikzpicture}\draw[black!35] (0,0) rectangle (\dimexpr\linewidth-.5pt\relax,#1);\end{tikzpicture}\par}
\newcommand{\Record}[2]{#1\par\Space{#2}}

% Scoped visual vocabulary: every state has the same clockwise lamp order.
\newcommand{\MRing}[5]{\begin{tikzpicture}[x=1in,y=1in,baseline=-.10in]
  \foreach \i in {1,...,4}{\coordinate (v\i) at ({90-90*(\i-1)}:#1);}
  \foreach \i/\j in {1/2,2/3,3/4,1/4}{
    \draw[black!65,line width=.8pt] (v\i)--(v\j);
    \ifnum#5=1\node[fill=white,inner sep=3pt,font=\normalsize] at ($(v\i)!.5!(v\j)$){\i\j};\fi}
  \foreach \i in {1,...,4}{\node[circle,draw=ink,line width=1pt,fill=white,minimum size=#2in,inner sep=0pt,font=#3\bfseries] at(v\i){\i};}
  \foreach \i in {#4}{\node[circle,draw=ink,line width=1pt,fill=ink,text=white,minimum size=#2in,inner sep=0pt,font=#3\bfseries] at(v\i){\i};}
\end{tikzpicture}}
\newcommand{\MMat}{\MRing{.90}{.56}{\large}{}{1}}
\newcommand{\MTarget}[1]{\MRing{.38}{.28}{\small}{#1}{0}}
\newcommand{\MRecord}{\MRing{.31}{.25}{\small}{}{0}}
\newcommand{\MTargetTask}[2]{\begin{minipage}[t]{3.28in}\centering
  #1\quad\MTarget{#2}\par\vspace{3pt}\raggedright
  Moves: \rule{2.20in}{.5pt}\par\vspace{8pt}
\end{minipage}}
\newcommand{\MTryTarget}[2]{\begin{minipage}[t]{3.28in}\centering
  #1\quad\MTarget{#2}\par\vspace{3pt}\raggedright
  Moves: \rule{2.20in}{.5pt}\par
  $\square$ Not yet\par\vspace{8pt}
\end{minipage}}
\newcommand{\MPair}[1]{\begin{tikzpicture}[x=1in,y=1in,baseline=-.12in]
  \draw[thick](0,0)--(1.12,0);
  \node[fill=white,inner sep=3pt,font=\small] at(.56,0){12};
  \foreach\i/\x in{1/0,2/1.12}{\node[circle,draw,fill=white,minimum size=.48in,inner sep=0pt,font=\large\bfseries] at(\x,0){\i};}
  \foreach\i/\x in{#1}{\node[circle,draw,fill=ink,text=white,minimum size=.48in,inner sep=0pt,font=\large\bfseries] at(\x,0){\i};}
\end{tikzpicture}}
\newcommand{\MTrial}[2]{\begin{minipage}{\linewidth}
  #1\quad\MPair{#2}\hfill$\xrightarrow{\text{press 12}}$\hfill\MPair{}\par
  \hspace*{.5in}ON before: \rule{.68in}{.5pt}\hfill ON after: \rule{.68in}{.5pt}\hspace*{.15in}\par
\end{minipage}\par\vspace{17pt}}
\newcommand{\MSortCard}[2]{\begin{minipage}[t]{1.64in}\centering
  #1\quad\MRing{.29}{.23}{\small}{#2}{0}\par
  ON: \rule{.45in}{.5pt}\par
  {\small made / not yet}\par\vspace{9pt}
\end{minipage}}
\newcommand{\MPath}[1]{\begin{tikzpicture}[x=1in,y=1in]
  \foreach\i in{1,...,4}{\coordinate(v\i)at({90-90*(\i-1)}:.57);}
  \foreach\i/\j in{1/2,2/3,3/4,1/4}{\draw[black!25,line width=.8pt](v\i)--(v\j);}
  \foreach\i/\j in{#1}{\draw[-{Stealth[length=7pt]},line width=2pt,shorten <=.18in,shorten >=.18in](v\i)--(v\j);}
  \foreach\i in{1,...,4}{\node[circle,draw,fill=white,minimum size=.35in,inner sep=0pt,font=\normalsize\bfseries]at(v\i){\i};}
\end{tikzpicture}}
\begin{document}
\fontsize{16}{20}\selectfont
\StudentHeader{Grades 2--3}
\Problem{1}{Put a two-sided counter on each big lamp, white side up: all OFF. Point to a line. Flip the counters at \emph{both} ends. Dark means ON.}
The line joining lamps 1 and 2 is called \textbf{12}. Replay these two moves on your big board.
\begin{center}
\begin{tabular}{c@{\hspace{.18in}}c@{\hspace{.18in}}c}
Start all OFF&Press 12&Then press 23\\[5pt]
\MTarget{}&\MTarget{1,2}&\MTarget{1,3}
\end{tabular}
\end{center}
Write the moves in order: \textbf{12, 23}.
\begin{center}\MMat\end{center}
Make the dark lamps white again by flipping pairs. Write your moves so a partner can replay them.
\vspace{10pt}\par Moves: \rule{5.85in}{.5pt}

\newpage
\StudentHeader{Grades 2--3}
\Problem{2}{Make each target picture on the big board. Reset all counters to white before each target. Flip both ends of one line at a time. Write your moves in order.}
\begin{center}\MMat\end{center}
\MTargetTask{A}{1,2}\hfill\MTargetTask{B}{1,3}\par
\vspace{11pt}
\MTargetTask{C}{2,4}\hfill\MTargetTask{D}{1,2,3,4}\par
\vspace{10pt}
Give a partner one of your move lists. Reset all OFF. Have your partner replay it and point to its target.

\newpage
\StudentHeader{Grades 2--3}
\Problem{3}{Try each target picture. Start all OFF each time. Flip both ends of one line at a time. Write a move list that works, or check ``not yet.''}
\begin{center}\MMat\end{center}
\MTryTarget{A}{1}\hfill\MTryTarget{B}{2,3}\par
\vspace{8pt}
\MTryTarget{C}{1,2,3}\hfill\MTryTarget{D}{1,4}\par
\vspace{6pt}
For a ``not yet'' picture, try a different move list on the board. Keep any list that works.

\newpage
\StudentHeader{Grades 2--3}
\Problem{4}{For each row, turn counters over by hand to set up its starting picture on the big board. Leave lamps 3 and 4 OFF. Then press 12 once. Darken the ON lamps in the answer picture. Count ON lamps before and after.}
\begin{center}\MMat\end{center}
\MTrial{A}{}
\MTrial{B}{1/0,2/1.12}
\MTrial{C}{1/0}

\newpage
\StudentHeader{Grades 2--3}
\Problem{5}{These are the targets you tried. For each picture, circle ``made'' or ``not yet.'' Count its ON lamps. Use the big board to try any picture again; reset all OFF each time.}
\vspace{10pt}
\MSortCard{A}{1,2}\hfill\MSortCard{B}{1,3}\hfill\MSortCard{C}{2,4}\hfill\MSortCard{D}{1,2,3,4}\par
\vspace{10pt}
\MSortCard{E}{1}\hfill\MSortCard{F}{2,3}\hfill\MSortCard{G}{1,2,3}\hfill\MSortCard{H}{1,4}\par
\begin{center}\MMat\end{center}
Compare the ON counts in the pictures you made.\par
Point to what is alike. Then compare the ``not yet'' pictures.

\newpage
\StudentHeader{Grades 2--3}
\Problem{6}{Follow each arrow path on the big board: flip both ends of each line, in order. Start all OFF for each path. Darken the ON lamps in its final picture.}
\vspace{7pt}
\begin{minipage}{3.4in}\centering
  A: \textbf{12, 23, 34}\par
  \MPath{1/2,2/3,3/4}
\end{minipage}\hfill
\begin{minipage}{2.8in}\centering
  Final picture:\par\MTarget{}
\end{minipage}\par
\vspace{9pt}
\begin{minipage}{3.4in}\centering
  B: \textbf{14, 43}\par
  \MPath{1/4,4/3}
\end{minipage}\hfill
\begin{minipage}{2.8in}\centering
  Final picture:\par\MTarget{}
\end{minipage}\par
\begin{center}\MMat\end{center}
Replay one path slowly. Point to each lamp that gets flipped twice. Point to each lamp flipped just once.

\newpage
\StudentHeader{Grades 2--3}
\Problem{7}{Start all OFF on the big board. Try some moves. Copy the result onto a little ring: darken its ON lamps. Reset and make different pictures. Include all OFF. Leave extra rings empty.}
\begin{center}\MMat\end{center}
\begin{center}
\begin{tabular}{@{}c@{\hspace{.67in}}c@{\hspace{.67in}}c@{\hspace{.67in}}c@{}}
\MRecord&\MRecord&\MRecord&\MRecord\\[-2pt]
$\square$ saved&$\square$ saved&$\square$ saved&$\square$ saved\\[16pt]
\MRecord&\MRecord&\MRecord&\MRecord\\[-2pt]
$\square$ saved&$\square$ saved&$\square$ saved&$\square$ saved\\[16pt]
\MRecord&\MRecord&\MRecord&\MRecord\\[-2pt]
$\square$ saved&$\square$ saved&$\square$ saved&$\square$ saved
\end{tabular}
\end{center}
Check ``saved'' under each picture you record, including all OFF. This keeps unused rings separate.

\newpage
\StudentHeader{Grades 2--3}
\Problem{8}{Look at the pictures you saved in Problem 7. Count the ON lamps in each. Write that count beside its little ring. Point to the pictures with the same count. Count how many different pictures you have.}
\vspace{9pt}
I saved \rule{.80in}{.5pt} different pictures.\par
\vspace{10pt}
Choose a picture whose moves you have not written down. Copy it here. Start all OFF on the big board and find a move list that makes it.
\begin{center}
\begin{minipage}{2.2in}\centering My target:\par\MTarget{}\end{minipage}
\end{center}
\begin{center}\MMat\end{center}
My moves: \rule{5.55in}{.5pt}\par
\vspace{16pt}\rule{\linewidth}{.5pt}\par
\vspace{8pt}
Reset all OFF and ask a partner to replay your list. Compare the result with your target.
\end{document}
